\section*{الفصل \ref{ch:g10:muscles-and-joints} --- العضلات والمفاصل}

\begin{solution}{exo:g10:muscles-and-joints:1}
\omterm{def:g10:muscles-and-joints:joint}{الغضروف} على نهايتي العظمين (انزلاق أملس، وتوزيع الحمل)؛ والمحفظة (تحكم
إغلاق \omterm{def:g10:muscles-and-joints:joint}{المفصل}) بسائلها الزليلي (التزليق)؛ \omterm{def:g10:muscles-and-joints:joint}{والأربطة} (تمسك العظمين معًا،
وتحد الحركات)؛ وفي الركبة، الهلالتان (توزعان الحمل).
\end{solution}

\begin{solution}{exo:g10:muscles-and-joints:2}
\omterm{def:g10:muscles-and-joints:muscle}{الوتر} يصل عضلة بعظم وينقل شدها؛ \omterm{def:g10:muscles-and-joints:joint}{والرباط} يصل عظمًا بعظم عبر \omterm{def:g10:muscles-and-joints:joint}{المفصل} ويحد
حركته. وما من واحد منهما ينقبض.
\end{solution}

\begin{solution}{exo:g10:muscles-and-joints:3}
العضلة لا تستطيع إلا الشد بالتقصر؛ ولا تستطيع الدفع. فلتحريك \omterm{def:g10:muscles-and-joints:joint}{مفصل} في
الاتجاهين تلزم قابضة وباسطة، كل منهما تنقض ما فعلته الأخرى.
\end{solution}

\begin{solution}{exo:g10:muscles-and-joints:4}
الالتواء: \omterm{def:g10:muscles-and-joints:joint}{رباط} تمدد أو تمزق. والتمزق العضلي: ألياف عضلية انقطعت.
والخلع: نهاية عظم دفعت خارج مفصلها، مع تمدد المحفظة \omterm{def:g10:muscles-and-joints:joint}{والأربطة} أو
تمزقها.
\end{solution}

\begin{solution}{exo:g10:muscles-and-joints:5}
الألياف البطيئة: بطيئة وضعيفة ولا تكل، غنية \omterm{def:g10:cells-common-unit:organelle}{بالمتقدرات} والشعيرات
الدموية، حمراء، تجري على التنفس. والألياف السريعة: سريعة وقوية، تتعب
في دقيقة، قليلة \omterm{def:g10:cells-common-unit:organelle}{المتقدرات}، شاحبة، تجري على مخزون الفوسفات \omterm{def:g10:cell-metabolism:fermentation}{والتخمر}.
\end{solution}

\begin{solution}{exo:g10:muscles-and-joints:6}
الحمل $2 \times 9.8 = \qty{19.6}{N}$؛ وقوة ذات الرأسين $19.6 \times 32/4
\approx \qty{157}{N}$.
\end{solution}

\begin{solution}{exo:g10:muscles-and-joints:7}
قد يأتي الدفع من أي من الاتجاهين، ولا تستطيع كل عضلة مقاومة إلا اتجاه
واحد. وانقباضهما معًا يصلّب \omterm{def:g10:muscles-and-joints:joint}{المفصل} في الاتجاهين دفعة واحدة؛ فمن أي
جهة جاء الدفع كانت إحداهما تشد ضده من قبل.
\end{solution}

\begin{solution}{exo:g10:muscles-and-joints:8}
يقصر الليف \qty{0.6}{cm}؛ \omterm{def:g10:muscles-and-joints:muscle}{والوتر} يحرك العظم \qty{0.6}{cm} على بعد
\qty{5}{cm} من المرفق، فتتحرك اليد
$0.6 \times 35/5 = \qty{4.2}{cm}$.
\end{solution}

\begin{solution}{exo:g10:muscles-and-joints:9}
الحركية تقتضي جوفًا ضحلًا وأربطة مرتخية؛ وتلك الرخاوة نفسها تتيح لرأس
العضد أن يغادر الجوف حين تدفعه قوة وراء المدى الذي تستطيع \omterm{def:g10:muscles-and-joints:joint}{الأربطة}
حفظه.
\end{solution}

\begin{solution}{exo:g10:muscles-and-joints:10}
التجفاف وفقد الملح بالتعرق، وتعب الألياف في نهاية السباق. والعلاج:
مطّ العضلة برفق واشرب ماءً بملح.
\end{solution}

\begin{solution}{exo:g10:muscles-and-joints:11}
فخذ العدّاء السريع 70\% ألياف سريعة والنسبة موروثة في معظمها؛ والتدريب
يضخم الألياف ويحسّن عدتها لكنه لا يحوّل أعدادها. أما غير المدرَّب فيبدأ
من النصف والنصف وعنده ألياف بطيئة أكثر لينميها.
\end{solution}

\begin{solution}{exo:g10:muscles-and-joints:12}
$4 \times 70 \times 9.8 \approx \qty{2.7}{kN}$. والإجهاد $2744/100
\approx \qty{27}{N/mm^2}$: أي نحو ربع إجهاد الانقطاع، بمعامل أمان
قدره أربعة.
\end{solution}

\begin{solution}{exo:g10:muscles-and-joints:13}
العضلة موصولة بالدم وصلًا غنيًا وفيها \omterm{def:g10:cells-common-unit:cell}{خلايا} احتياطية تعيد بناء
الألياف: فأسابيع. \omterm{def:g10:muscles-and-joints:joint}{والأربطة} أوعيتها قليلة، فتصل المغذيات \omterm{def:g10:cells-common-unit:cell}{وخلايا} الإصلاح
ببطء: فأشهر. أما \omterm{def:g10:muscles-and-joints:joint}{الغضروف} فلا أوعية دموية فيه البتة ولا تكاد تكون فيه
\omterm{def:g10:cells-common-unit:cell}{خلايا} منقسمة؛ ويغذيه السائل الزليلي وحده ولا يكاد يصلح.
\end{solution}

\begin{solution}{exo:g10:muscles-and-joints:14}
كمال الأجسام يغلّظ الألياف (لييفات أكثر في كل ليف) والتمطيط يطيل وحدة
العضلة \omterm{def:g10:muscles-and-joints:muscle}{والوتر} ويوسع احتمال المحفظة. وعدد الألياف يثبت باكرًا في
الحياة؛ والعضلة البالغة تنمو أو تضمر بتغير حجم أليافها الموجودة.
\end{solution}

\begin{solution}{exo:g10:muscles-and-joints:15}
على السلالم يحمل \omterm{def:g10:muscles-and-joints:joint}{المفصل} عدة أمثال وزن الجسم؛ \omterm{def:g10:muscles-and-joints:joint}{والغضروف} الأنحف ينقل ذلك
الضغط إلى العظم تحته، وفيه أعصاب فيؤلم. ومع \omterm{def:g10:muscles-and-joints:joint}{غضروف} أقل تصير الهلالتان
آخر ما يوزع الحمل. والعضلات القوية حول الركبة تمتص جزءًا من الصدمة
وتثبّت \omterm{def:g10:muscles-and-joints:joint}{المفصل}، فتخفض الضغط الأعظمي على ما بقي من \omterm{def:g10:muscles-and-joints:joint}{الغضروف}.
\end{solution}

\begin{solution}{pb:g10:muscles-and-joints:1}
\textbf{1.} $70 \times 9.8 = \qty{686}{N}$.

\textbf{2.} $686 \times 10/5 \approx \qty{1.4}{kN}$.

\textbf{3.} $686 \times 20/5 \approx \qty{2.7}{kN}$.

\textbf{4.} مثلا وزن الجسم وأربعة أمثاله: ففي كل درجة تشد رباعية
الرؤوس بوزن أربعة أشخاص، ولذلك تحترق الفخذان في صعود طويل.

\textbf{5.} أوتار المأبض، في خلف الفخذ؛ وهي ترتخي وتتمدد بينما تبسط
رباعية الرؤوس الركبة، وتنقبض قليلًا لتثبيت \omterm{def:g10:muscles-and-joints:joint}{المفصل}.

\textbf{6.} قوة الأرض $3 \times 686 \approx \qty{2.06}{kN}$؛ وقوة
\omterm{def:g10:muscles-and-joints:muscle}{الوتر} $2058 \times 20/5 \approx \qty{8.2}{kN}$، أي اثني عشر ضعف وزن
الجسم.

\textbf{7.} $8232/120 \approx \qty{69}{N/mm^2}$: أي نحو ثلثي إجهاد
الانقطاع. فكل هبوط يستهلك معظم هامش \omterm{def:g10:muscles-and-joints:muscle}{الوتر}.

\textbf{8.} الحمل هو $8.2 + 2.1 \approx \qty{10.3}{kN}$. ومع
الهلالتين يكون $10300/600 \approx \qty{17}{N/mm^2}$؛ ومن دونهما
يكون $10300/100 \approx \qty{103}{N/mm^2}$، أي ستة أمثال.

\textbf{9.} ستة أمثال الضغط على \omterm{def:g10:muscles-and-joints:joint}{الغضروف} نفسه في كل هبوط: فيسحق السطح
أسرع مما يستطيع أن يصان، فيبلى.

\textbf{10.} غشاء من السائل بين الغضروفين، لينزلقا أثناء الانثناء بلا
احتكاك جاف، ومع تخميد طفيف للصدمة.

\textbf{11.} التواء الظنبوب تحت عظم الفخذ، مع انثناء جانبي: أي دورانات
لم تبن لها الركبة، \omterm{def:g10:muscles-and-joints:joint}{والأربطة} موجودة لمقاومتها.

\textbf{12.} رباعية الرؤوس لا تشد إلا على طول الرجل، فتبسط الركبة؛ ولا
فعل لها ضد الالتواء. ثم إن الإصابة تقع في أجزاء من مئة من الثانية،
أسرع من أي استجابة عضلية.

\textbf{13.} دم من أوعية \omterm{def:g10:muscles-and-joints:joint}{الرباط} الممزق والمحفظة، مع سائل زليلي زائد:
فقد خرقت المحفظة أو التهبت وامتلأ \omterm{def:g10:muscles-and-joints:joint}{المفصل}.

\textbf{14.} \omterm{def:g10:muscles-and-joints:joint}{الرباط}، وإمداده الدموي ضعيف، يحتاج إلى أشهر وإلى جراحة
غالبًا؛ أما الهلالة الممزقة، وهي \omterm{def:g10:muscles-and-joints:joint}{غضروف} يكاد يخلو من الأوعية، فلا تشفى
عادةً ويشذَّب الجزء الممزق أو يخاط.

\textbf{15.} أوتار المأبض مضادة رباعية الرؤوس وهي تشد الظنبوب إلى
الخلف؛ وزوج قوي ينقبض معًا يصلّب الركبة ضد الانزلاق الأمامي الذي لم
يعد \omterm{def:g10:muscles-and-joints:joint}{الرباط} الممزق يمنعه، فيحمي \omterm{def:g10:muscles-and-joints:joint}{المفصل} قبل الجراحة وبعدها.

\textbf{16.} نحو 4000 مرة في اليوم عند \qty{1.4}{kN} أو أكثر.

\textbf{17.} نحو $\num{1.5e6}$ تحميلة في السنة. والتهاب \omterm{def:g10:muscles-and-joints:muscle}{الوتر} أذى
يتراكم أسرع من الإصلاح البطيء لنسيج \omterm{def:g10:muscles-and-joints:muscle}{الوتر}؛ ولا شيء يتيح للإصلاح أن
يلحق إلا الراحة، لأن القوة لا تخفض حمل كل خطوة.

\textbf{18.} الهبوطات $40 \times 8.2 \approx \qty{330}{kN}$ إجمالًا؛
والخطوات $4000 \times 1.4 \approx \qty{5600}{kN}$ --- فالخطوات تحمل
سبعة عشر ضعف الحمل الكلي، لكن الهبوطات تحمل الذرى التي تهدد بالانقطاع.

\textbf{19.} مضاعفة قوة الأرض تضاعف قوتي \omterm{def:g10:muscles-and-joints:muscle}{الوتر} \omterm{def:g10:muscles-and-joints:joint}{والمفصل} في كل خطوة؛
وركوب الدراجة والسباحة يشغلان العضلات مع تحميل الركبة دون وزن الجسم
بكثير.

\textbf{20.} يحمل \omterm{def:g10:muscles-and-joints:muscle}{وتر} رباعية الرؤوس أربعة أمثال وزن الجسم على درجة
ونحو اثني عشر عند هبوط؛ أما \omterm{def:g10:muscles-and-joints:joint}{الأربطة}، وهي ما لا تختبره تلك الأحمال
قط، فكانت النسيج الذي وجدته الاستدارة.
\end{solution}
